\documentclass[12pt]{article}
\usepackage[utf8]{inputenc}
\usepackage[margin=0.75in]{geometry}
\usepackage{amsmath}\usepackage{amssymb}\usepackage{enumitem}\usepackage{multicol}\usepackage{tcolorbox}
\setlength{\parindent}{0pt}
\begin{document}

\begin{center}{\Large \textbf{Describe transformations of parent functions from the equation}}\\[2pt]{\small Pre-Cal \textbar\ CK-12: Graphical Transformations}\end{center}\vspace{0.25cm}

{\large\textbf{Key Idea}}\par\smallskip
Transformations move or reshape a parent graph. In $y=a\,f(x-h)+k$: $h$ shifts left/right, $k$ shifts up/down, $a$ stretches (and flips if negative).\par


{\large\textbf{Key Terms}}\par\smallskip
\begin{itemize}[leftmargin=1.4em]
  \item \textbf{Translation} — a slide (from $h$ or $k$)
  \item \textbf{Reflection} — a flip (from a negative $a$)
\end{itemize}


{\large\textbf{Worked steps}}\par\smallskip
Step 1: Match the equation to $a\,f(x-h)+k$.\par
Step 2: Read $h$ (opposite sign), $k$, and $a$.\par
Example: $y=-(x-2)^2+5$ — right 2, up 5, and flipped over the x-axis.\par


{\large\textbf{You Try}}\par\smallskip
Describe $y=(x+3)^2-1$ from the parent $y=x^2$.\par

\end{document}